\documentclass[11pt]{article} %fonttikoko
\usepackage{graphicx} % Required for inserting images
\usepackage[hidelinks]{hyperref} % linkit sisälllysluettelossa
\usepackage[left=1cm,right=1cm,bottom=1in]{geometry}
%\usepackage{verse} % sisältyy laulupankki -pakettiin
\usepackage{multicol}


\usepackage{laulupankki}

%-----------------------------------------------------------------------------
%
%       MEIDÄN HIENO PAKETTIMUOTOINEN LAULUPANKKI laulupankki.sty
%
%-----------------------------------------------------------------------------



\usepackage{titlesec} % Otsikoiden formatointia varten

\usepackage{blindtext}
\usepackage[finnish]{babel}

\usepackage{tocloft} % Formationtia varten

\renewcommand{\cftsecleader}{\cftdotfill{\cftdotsep}}   %   Piste-erottimet sisällysluetteloon

\setlength{\cftbeforesecskip}{2pt}   % tiivistys sisällysluetteloon section
\setlength{\cftbeforesubsecskip}{1pt}   % tiivistys sisällysluetteloon subsection


%
%           Seuraavalla komennolla lisätään tilaa sisällysluetteloon
%           Laulun numeron ja otsikon väliin.
%

\setlength{\cftsecnumwidth}{20pt} % väli sisällysluettelon numeroinnin ja otsikoinnin välillä


%\renewcommand{\thesubsection}{\thesection67.} %custom pääte alaotsikoihin


%
%       Määritellään laulujen otsikot pienemmiksi, ja lisätään piste numeron perään
%

\titleformat{\section}
  {\normalfont\bfseries} %fonttikoko normal, bold 
  {\thesection.}
  {1em}
  {}

%
%       Sama operaatio alaotsikoille, sikäli mikäli käytät alaotsikoita esim.
%       vaihtoehtoisia sanoituksia varten laulupankissa
%


\titleformat{\subsection}
  {\normalfont\bfseries} %fontti normal, bold 
  {\thesubsection}
  {1em}
  {}

%
%           Määritellään seuraavaksi laulujen otsikoitden asettelu, koska asettelu
%           menee vituiksi otsikon formatoinnin yhteydessä
%

\titlespacing*{\section}
{1,2cm} %tilaa vasemmalla
{2pt} %tilaa päällä
{0,2pt} %tilaa alla


%\renewcommand\thesubsection{\thesection\arabic{subsection}.} % pisteet alaotsikkonumeroiden perään

\titlespacing*{\subsection} %samat alaotsikoille
{1,2cm}
{2pt}
{0,2pt}


%\renewcommand{\thepage}{67} %custom sivunumero 67-sitseillä :D

%
%           Voit siirrellä lauluja sivulta toiselle komentoyhdistelmällä
%           \vfill\null \columnbreak
%           Tämä kuitenkin on työlästä, joten määritellään uusi komento tätä varten
%

\newcommand{\cb}{\vfill\null \columnbreak}

%           
%           Nyt voit komennolla \cb siirtää seuraavan laulun alun seuraavalle
%           kolumnille tai sivulle.
%


%---------------------------------
%       ETUSIVU
%---------------------------------

%\sffamily
\begin{document}

\thispagestyle{empty}
\vspace*{1cm}
\begin{center}
%\includegraphics[width=0.9\linewidth]{kuva.png}\\
\vspace{1cm}
\LARGE \textbf{Laululäsy}
\end{center}

\clearpage

%---------------------------------
%       SÄÄNNÖT
%---------------------------------

\begin{center}
\textbf{Sitsietiketti}
\vspace{0.3cm}



\begin{itemize}
    \item Seremoniamestarit pyytävät puheenvuoron sauvan avulla.
    \begin{itemize}
        \item Yksi kopautus: Seremoniamestarit pyytävät hiljaisuutta
        \item Kaksi kopautusta: Seramoniamestarit pyytävät hiljaisuutta JA heillä on asiaa
        \item Kolme kopautusta: Seremoniamestarit pyytävät hiljaisuutta JA heillä on asiaa JA heillä on jano
    \end{itemize}
    \item Jokaisen puheenvuoron ja laulun lopuksi skoolataan. Skål-järjestys: etuvasen, etuoikea, eteen
    \item Kaikille juomille tulee laulaa kerran ennen niiden avaamista. Kun juoma on avattu, sitä saa juoda milloin vain.
    \item Vessassa saa käydä sitsien aikana vessahattu päässä. Muuten pöydästä ei saa nousta kuin vain tauolla. Vältetään kuitenkin tilassa kulkemista laulujen ja puheiden aikana.
    \item Hyvään sitsietikettiin kuuluu hyvä käytös. Ethän pilaa muiden sitsaajien iltaa omalla toiminnallasi.
    \item Puhelimen käyttäminen sitsien aikana on kielletty.
\end{itemize}

\textbf{Puheenvuorot}
\vspace{0.3cm}
\begin{itemize}
    \item Jokainen sitseillä olija, voi pitää puheenvuoron. Puheenvuoron saa nousemalla ylös ja sanomalla kovaan ääneen \textit{"Kilinkolin, kilinkolin"}.
    \item Esittele itsesi puheenvuoron aluksi. Esittelyn jälkeen voit pitää puheenvuoron. Puheenvuoron lopuksi saat halutessasi toivoa laulua.
    \item Lauluja toivotaan vain laululäsystä. Lisäsäkeistöjen laulaminen on sallittua.
    \item Annetaan mahdollisimman monelle mahdollisuus pitää puheenvuoroja. Ethän siis pyydä puheenvuoroa jatkuvasti itsellesi ;)
\end{itemize}

\textbf{Välihuutelut}
\vspace{0.3cm}
\begin{itemize}
    \item \textit{Skål!} Huudetaan yhdessä jokaisen laulun lopuksi ja skoolataan sitsietiketin mukaan.
    \item \textit{Mellis!} Suomennos sanoista mellan sup, eli välihörppy. Mellis huudetaan kesken laulun, jolloin laulaminen lopetetaan välihörpyn ajaksi. Hörpyn jälkeen laulua jatketaan.
    \item \textit{Omstart!} Laulu tai säkeistö aloitetaan alusta. Huudetaan usein jos laulaminen ei suju kuten sen kuuluisi sujua.
    \item \textit{Tempo!} Laulun aikana kuuluva huudahdus, mikä tarkoittaa että laulun tempoa nostetaan eli laulaminen tapahtuu nopeampaa.
\end{itemize}
\end{center}


\clearpage

%---------------------------------
%       SISÄLLYSLUETTELO
%---------------------------------

\begin{multicols}{2}

\tableofcontents

\end{multicols}


\clearpage


%---------------------------------
%           
%       ITSE BIISIT ALKAA
%
%---------------------------------



\begin{multicols}{2}

    \helangar
    
    %---------------------------------
    %           vesi
    %---------------------------------
    \aquavera

    \cb

    \vesihenkilo

    %---------------------------------
    %           kalja
    %---------------------------------
    \nautintoainekaanon
    \kaljaa
    


   % \ystavasasnapsien
    
    %---------------------------------
    %           siideri
    %---------------------------------
    \sensiiderin

    %---------------------------------
    %           lonkero
    %---------------------------------
    \lonkero


    %---------------------------------
    %           vahvat
    %---------------------------------
   % \vauhtijuoma
    \viinamaa

    \koskenkorvaa
    

    %---------------------------------
    %           jallu
    %---------------------------------
    \jallutahdenalla

    
    %---------------------------------
    %           viini
    %---------------------------------
    \skumppaapunkunkaa
    %\boheemirapsakka

    %---------------------------------
    %           muita juomia
    %---------------------------------
    %\gambinacocktailjuoma


    %---------------------------------
    %           RYYBS+KRAPS+KRAPULA
    %---------------------------------
    
    \karhuviina
    %\aamullavarhain
    \laulutaidottomille
    \minnet
    \ikuinenkanni

    \tanaanotetaan
    %\larvit
    %\eurovision
    %\aidinpienetsienet



    


    %---------------------------------
    %           leikkilauluja
    %---------------------------------
    \testamentti
    \pikkukakkonen
    \varpunen


    %---------------------------------
    %           matfys lauluja
    %---------------------------------
    
    %\laser
    %\osittaisderivaattalaulu
    \derivaattanollanarvoinen

    %---------------------------------
    %           muuta hassua
    %---------------------------------
    %\eukkosi
    \muinaistenvalssi

    \mandariinimies
    %\setoinenmandariinilaulu
    \mandariinihunningolla
    
    \volga
    
    \miesulossinusta
    %\sauron
    
    \vanhustenpoyta
    \amorinfinitus
    \onomato
    %\sahtilaulu
    
    %---------------------------------
    %           kauneimmat oululaulut
    %---------------------------------
    \toppila
    \oulumaa
    \oululainenbaarikierros
    %\tuuppatanneouluun
    \kukakakkaaaltaaseen
    %\kopille
    %\kopinraja
    %\paljukansa
    \kaijonvalot
    \ukkometso

    %---------------------------------
    %           lisää hauskoja lauluja
    %---------------------------------
    \sitsipoytahajotettu
    %\patkismunkki
    %\hulivilipojanhumppa
    
    %---------------------------------
    %           vittu
    %---------------------------------
    %\vlaulu
    \vitutuslaulu
    \viskadrickahelanatten
    
    %---------------------------------
    %           meemi bängereitä
    %---------------------------------

    %\grandiosa
    \yollataas
    \siljalinespecial
    %\sandstorm
    \piilossaviinat 
    \exculaulu
    \sornaigusha
    %\hittiputki
    \annannen
    %\schro

    %---------------------------------
    %           loppu- ja välilaulut
    %---------------------------------
    \henkilokunta 
    \fuksilaulu

    

\end{multicols}

\clearpage
\begin{center}
    

\vspace{1cm}
\textbf{Turvallisemman tilan periaatteet}\\
Tapahtumassa noudatetaan --- Turvallisemman tilan periaatteita\\

\vspace{1.0cm}
\textbf{Tämän tapahtuman häirintäyhteyshenkilöt:}\\

\vspace{1cm}
\textbf{--- häirintäyhteyshenkilöt tavoittaa osoitteesta}\\

\end{center}
\end{document}